% Generated draft: requires mathematical and source review.
\documentclass[11pt]{article}
\usepackage[margin=25mm]{geometry}
\usepackage{fontspec}
\setmainfont{Latin Modern Roman}
\usepackage{amsmath,amssymb,mathtools,unicode-math}
\setmathfont{Latin Modern Math}
\usepackage{enumitem,needspace,xurl,hyperref}
\setlength{\parindent}{0pt}
\setlength{\parskip}{5pt}
\setlength{\emergencystretch}{3em}
\newcommand{\N}{\mathbb{N}}
\newcommand{\Z}{\mathbb{Z}}
\newcommand{\Q}{\mathbb{Q}}
\newcommand{\R}{\mathbb{R}}
\newcommand{\C}{\mathbb{C}}
\newcommand{\F}{\mathbb{F}}
\newcommand{\A}{\mathbb{A}}
\newcommand{\PP}{\mathbb{P}}
\newcommand{\E}{\mathbb{E}}
\newcommand{\eps}{\varepsilon}
\newcommand{\dd}{\,\mathrm{d}}
\begin{document}
\section*{UnsolvedMath ID: 30006385}
\textbf{Source problem code:} catalogue entry
\subsection{Definitions and mathematical statement}
Let \(n\) denote the block length and let \(\F_q\) be a finite field, so that \(q\) is a prime power. A linear code is a linear subspace \(C_n\leq \F_q^n\). For a subset \(B\subseteq\F_q^n\), define its list size in the code by
\[
L(C_n,B):=|C_n\cap B|.
\]
The supplied problem asks whether a random linear code has, with high probability, list size at most \(q^{\varepsilon n}(1+o(1))\) simultaneously in every admissible ball defined by constraints beyond coordinate-wise restrictions.

The source record does not define these balls, specify the distribution or dimension of the random code, or give the admissible values and dependencies of the parameters. Consequently, it does not determine a complete mathematical conjecture. To display the intended bound without supplying those missing hypotheses, write \(\mathcal B_n\subseteq\mathcal P(\F_q^n)\) for the unspecified family of admissible balls, where \(\mathcal P\) denotes the power set. Under the usual interpretation of a uniform multiplicative \(o(1)\) term, the requested assertion would have the schematic form
\[
\Pr\!\left[
  \forall B\in\mathcal B_n,\quad
  L(C_n,B)\leq q^{\varepsilon n}(1+\delta_n)
\right]\longrightarrow 1
\qquad(n\longrightarrow\infty),
\]
for some deterministic nonnegative sequence \(\delta_n\to0\). Here the probability is over the unspecified random-code ensemble. The field size, code dimension or rate, ball parameters, range of \(\varepsilon\), and permitted dependencies of \(\delta_n\) remain unresolved. In particular, the displayed formula is a conditional template for the supplied question, not a fully specified conjecture. The phrase ``every such ball'' requires missing antecedent context; no metric-ball definition or alternative constraint family is inferred.
\subsection{Short English statement}
The question asks whether every admissible ball contains at most \(q^{\varepsilon n}(1+o(1))\) codewords of a random linear code, with probability tending to one. A precise conjecture requires the missing ball definition, random-code distribution, code rate, and parameter hypotheses. The supplied record labels the problem open; this draft does not independently verify that status.
\subsection{Sources}
\begin{itemize}
\item Local registry: \texttt{data/unsolvedmath/problems.json}, with \texttt{catalog/problems.json} as the supplemental registry; supplied permanent UnsolvedMath ID \(30006385\), source problem code \texttt{catalogue entry}, title ``List Sizes Beyond Coordinate-Wise Constraints.'' This attribution uses the supplied record; the registry files were not independently inspected.
\item Supplied report citation: \emph{New Mathematical Directions in Coding Theory (2026)}, Oberwolfach Reports, DOI \texttt{10.4171/owr/2025/41}, pp.\ 2187--2242. \url{https://doi.org/10.4171/owr/2025/41}. The record attributes the extracted question to an open-problem discussion in this report and supplies proposed year 2025, but names no individual proposer. The citation's 2026 title date and its 2025 DOI/report designation are retained as supplied.
\item Supplied literature source: ``Oberwolfach Report 41 2025.'' \url{https://ems.press/content/serial-article-files/52245}.
\item Supplied literature source: \emph{Near-Optimal List-Recovery of Linear Code Families}. \url{https://arxiv.org/abs/2502.13877}.
\item Literature assessment dated \texttt{2026-08-22}: supplied dated curation, not independently verified. The record labels the problem open and reports that the Oberwolfach survey leaves this above-capacity question open, that a violating list is exponential in \(n\) and beyond the current local-property framework, and that later random-linear list-recovery results cover important capacity regimes but not the stated uniform bound. These are supplied assessments requiring independent verification of time-sensitive status and cited claims.
\item The record labels its statement \texttt{corrected\_verified} and states that the official report was checked and notation repaired. These are supplied source-verification claims, not verification performed for this draft. No report page or numbered passage locating this particular question is supplied beyond the report-wide page range.
\end{itemize}
\end{document}
